\section{Differentiering}

\subsection{Partielt Afledte}
En partielt afledt af en funktion $f(x,y)$ behandler den ene som konstant og differentierer i forhold til den anden.
\eksempel{
    \(
        \frac{\partial f}{\partial x} (3x^2y + 2y^3) = 
        \frac{\partial f}{\partial x} (3x^2y) =
        3y \cdot \frac{\partial f}{\partial x} (x^2) =
        3y \cdot 2x =
        6xy
    \)
}

\subsubsection{Gradienter}
Gradienten af en funktion $f(\vec x)$ er en vektor, der indeholder de partielt afledte i forhold til hver variabel:
\[
    \nabla f(x,y) = \left( \frac{\partial f}{\partial x}, \frac{\partial f}{\partial y} \right)
\]

\subsection{Hesse Matricer}
Hesse matricen af en funktion $f(\vec x)$ er en kvadratisk matrix, der indeholder de andenordens partielt afledte:
\begingroup
\renewcommand*{\arraystretch}{1.5}
\[
    H_f(x,y) = \begin{bmatrix}
        \frac{\partial^2 f}{\partial x^2} & \frac{\partial^2 f}{\partial x \partial y} \\
        \frac{\partial^2 f}{\partial y \partial x} & \frac{\partial^2 f}{\partial y^2}
    \end{bmatrix}
    \qquad
    H_f(x,y,\ldots,z) = \begin{bmatrix}
        \frac{\partial^2 f}{\partial x^2} & \frac{\partial^2 f}{\partial x \partial y} & \cdots & \frac{\partial^2 f}{\partial x \partial z} \\
        \frac{\partial^2 f}{\partial y \partial x} & \frac{\partial^2 f}{\partial y^2} & \cdots & \frac{\partial^2 f}{\partial y \partial z} \\
        \vdots & \vdots & \ddots & \vdots \\
        \frac{\partial^2 f}{\partial z \partial x} & \frac{\partial^2 f}{\partial z \partial y} & \cdots & \frac{\partial^2 f}{\partial z^2}
    \end{bmatrix}
\]
\endgroup

\bemærkBig{
    \( \frac{\partial^2 f}{\partial y \partial x} = \frac{\partial^2 f}{\partial x \partial y} \) for kontinuerete funktioner.
    \kilde{Clairaut's sætning}

    $\implies$ Hesse matricen af en kontinuert funktion er symmetrisk.
}

\subsection{Jacobi Matricer}
Givet vektorfunktion $\vec f: \mathbb{R}^n \to \mathbb{R}^m$, er Jacobi matricen $J_{\vec f}$ en $m \times n$ matrix, der indeholder de førsteordens partielt afledte:
\begingroup
\renewcommand*{\arraystretch}{1.5}
\[
    J_{\vec f}(\vec x) = \begin{bmatrix}
        \frac{\partial f_1}{\partial x_1} & \frac{\partial f_1}{\partial x_2} & \cdots & \frac{\partial f_1}{\partial x_n} \\
        \frac{\partial f_2}{\partial x_1} & \frac{\partial f_2}{\partial x_2} & \cdots & \frac{\partial f_2}{\partial x_n} \\
        \vdots & \vdots & \ddots & \vdots \\
        \frac{\partial f_m}{\partial x_1} & \frac{\partial f_m}{\partial x_2} & \cdots & \frac{\partial f_m}{\partial x_n}
    \end{bmatrix}
\]
\endgroup

\subsubsection{Jacobianten}
Jacobianten er determinanten af Jacobi matricen.

\subsection{Differentiabilitet}
En funktion \emph{fra en åben mængde} er differentiabel i $x_0$ hvis der eksisterer
$c \in \mathbb{R}$ og en Epsilon-funktion $\varepsilon(h)$ sådan at den lineære afbildning:

\[
    f(x_0+h) - f(x_0) = c \cdot h + |h| \cdot \varepsilon(h) 
\]

for alle $h \in \mathbb{R}$.

\kilde{Lemma 3.6.1}

\subsubsection{Epsilon-funktioner}
En Epsilon-funktion $\varepsilon(h)$ er en vilkårlig funktion der går mod 0 når $h$ går mod 0.